% --- Archon physics formalization source begin ---
% archon:physics
% archon:covers PhyXMiniProblems/problem_phyx_mini_0597.lean
% archon:source-report reports/phyx_mini/problem_phyx_mini_0597.source.json

\chapter{Physics problem phyx\_mini\_0597}
\label{ch:PhyXMiniProblems_problem_phyx_mini_0597}

\paragraph{Problem source.}
\#\# Physical scenario

Figure gives the velocity \textbackslash{}( u' \textbackslash{}) of the particle relative to frame \textbackslash{}( S' \textbackslash{}) for a range of values for \textbackslash{}( v \textbackslash{}). The vertical axis scale is set by \textbackslash{}( u'\_a = 0.800c \textbackslash{}).

\#\# Question

What value will \textbackslash{}( u' \textbackslash{}) have if \textbackslash{}( v \textbackslash{}rightarrow c \textbackslash{})?

\#\# Answer choices

- A: A: 0

- B: B: c

- C: C: -c

- D: D: 0.80c

\#\# Figure caption (auxiliary; use the image as primary evidence)

The image is a graph with a line plot. The x-axis is labeled with "0," "0.2c," and "0.4c," while the y-axis is labeled with "u'." There is also an additional label on the y-axis, "u'a." The graph shows a decreasing curve from left to right. The grid is present in the background, and the line appears to be bold.

\#\# Dataset metadata

Relativity; Physical Model Grounding Reasoning, Predictive Reasoning

\paragraph{Recorded answer/context.}
C: C: -c

\paragraph{Figure/image path.}
/root/proposal\_for\_physic/science-mango/phyx\_mini\_run/phyx\_data/test\_image/597.png

\paragraph{Formalization target.}
create a compiling Lean file with sorry bodies at `PhyXMiniProblems/problem\_phyx\_mini\_0597.lean`. The Lean declarations must preserve the physical quantities, dimensions or dimensional roles, figure labels, governing-law hypotheses, and final relation expressed by this problem.
Use Mathlib/Physlib names found through LeanExplore where available. If a physics API is missing, introduce faithful local abstractions rather than scalar placeholder aliases.

\begin{theorem}[Physics formalization target]
\label{thm:physics:phyx_mini_0597:target}
\lean{PhyXMiniProblems.ProblemPhyXMini0597.velocity_in_moving_frame_tends_to_negative_c}
\uses{def:physics:phyx-mini-0597:phyxminiproblems-problemphyxmini0597-relativisticvelocitygraphsetup, def:physics:phyx-mini-0597:phyxminiproblems-problemphyxmini0597-particlelongitudinalvelocityinmeterspersecond, def:physics:phyx-mini-0597:phyxminiproblems-problemphyxmini0597-matchessuppliedvelocitygraph, def:physics:phyx-mini-0597:phyxminiproblems-problemphyxmini0597-hasphysicalspeedoflight, def:physics:phyx-mini-0597:phyxminiproblems-problemphyxmini0597-satisfiesonedimensionallorentzvelocitylaw}
The assigned autoformalize agent should translate this physics problem into Lean declarations in the covered file, with theorem and lemma proof bodies written as `by sorry`.
\end{theorem}
\begin{proof}
This is an autoformalization task, not a proof task. Produce statements that can later be proved by the physics prover without weakening the physical model.
\end{proof}

% --- Archon declaration topology begin ---
\section{Declaration topology}
\begin{definition}[Inertial Frame Label]
  \label{def:physics:phyx-mini-0597:phyxminiproblems-problemphyxmini0597-inertialframelabel}
  \lean{PhyXMiniProblems.ProblemPhyXMini0597.InertialFrameLabel}
  \uses{def:physics:phyx-mini-0597:phyxminiproblems-problemphyxmini0597-relativespeedinmeterspersecond}
  The laboratory frame S and the collinear moving frame S'(v)
\end{definition}
\begin{proof}
  This is the declared model definition of Inertial Frame Label.
\end{proof}
\begin{definition}[relative Speed In Meters Per Second]
  \label{def:physics:phyx-mini-0597:phyxminiproblems-problemphyxmini0597-relativespeedinmeterspersecond}
  \lean{PhyXMiniProblems.ProblemPhyXMini0597.relativeSpeedInMetersPerSecond}
  \uses{def:physics:phyx-mini-0597:phyxminiproblems-problemphyxmini0597-inertialframelabel}
  The relative speed assigned to each named frame, measured from S
\end{definition}
\begin{proof}
  This is the declared model definition of relative Speed In Meters Per Second.
\end{proof}
\begin{definition}[longitudinal Velocity Fraction Of C]
  \label{def:physics:phyx-mini-0597:phyxminiproblems-problemphyxmini0597-longitudinalvelocityfractionofc}
  \lean{PhyXMiniProblems.ProblemPhyXMini0597.longitudinalVelocityFractionOfC}
  The longitudinal three-velocity component as a dimensionless fraction of the vacuum speed of light. This is a scalar projection of a physical Physlib four-velocity, not a scalar replacement for that velocity
\end{definition}
\begin{proof}
  This is the declared model definition of longitudinal Velocity Fraction Of C.
\end{proof}
\begin{definition}[Horizontal Axis Tick Label]
  \label{def:physics:phyx-mini-0597:phyxminiproblems-problemphyxmini0597-horizontalaxisticklabel}
  \lean{PhyXMiniProblems.ProblemPhyXMini0597.HorizontalAxisTickLabel}
  Labels printed along the horizontal v-axis
\end{definition}
\begin{proof}
  This is the declared model definition of Horizontal Axis Tick Label.
\end{proof}
\begin{definition}[Vertical Axis Label]
  \label{def:physics:phyx-mini-0597:phyxminiproblems-problemphyxmini0597-verticalaxislabel}
  \lean{PhyXMiniProblems.ProblemPhyXMini0597.VerticalAxisLabel}
  Labels printed along the vertical velocity axis
\end{definition}
\begin{proof}
  This is the declared model definition of Vertical Axis Label.
\end{proof}
\begin{definition}[Supplied Velocity Graph]
  \label{def:physics:phyx-mini-0597:phyxminiproblems-problemphyxmini0597-suppliedvelocitygraph}
  \lean{PhyXMiniProblems.ProblemPhyXMini0597.SuppliedVelocityGraph}
  \uses{def:physics:phyx-mini-0597:phyxminiproblems-problemphyxmini0597-horizontalaxisticklabel, def:physics:phyx-mini-0597:phyxminiproblems-problemphyxmini0597-verticalaxislabel}
  Typed data carried by image 597. The plotted ordinate and tick positions are dimensional scalar readouts in metres per second. Their numerical calibration and physical interpretation are stated separately as evidence
\end{definition}
\begin{proof}
  This is the declared model definition of Supplied Velocity Graph.
\end{proof}
\begin{definition}[Relativistic Velocity Graph Setup]
  \label{def:physics:phyx-mini-0597:phyxminiproblems-problemphyxmini0597-relativisticvelocitygraphsetup}
  \lean{PhyXMiniProblems.ProblemPhyXMini0597.RelativisticVelocityGraphSetup}
  \uses{def:physics:phyx-mini-0597:phyxminiproblems-problemphyxmini0597-inertialframelabel, def:physics:phyx-mini-0597:phyxminiproblems-problemphyxmini0597-suppliedvelocitygraph}
  The independent physical setup. A single particle has a genuine Physlib four-velocity in each frame, while speedOfLightInMetersPerSecond supplies the dimensional conversion between fractions of c and graph readouts
\end{definition}
\begin{proof}
  This is the declared model definition of Relativistic Velocity Graph Setup.
\end{proof}
\begin{definition}[particle Longitudinal Velocity In Meters Per Second]
  \label{def:physics:phyx-mini-0597:phyxminiproblems-problemphyxmini0597-particlelongitudinalvelocityinmeterspersecond}
  \lean{PhyXMiniProblems.ProblemPhyXMini0597.particleLongitudinalVelocityInMetersPerSecond}
  \uses{def:physics:phyx-mini-0597:phyxminiproblems-problemphyxmini0597-inertialframelabel, def:physics:phyx-mini-0597:phyxminiproblems-problemphyxmini0597-longitudinalvelocityfractionofc, def:physics:phyx-mini-0597:phyxminiproblems-problemphyxmini0597-relativisticvelocitygraphsetup}
  The particle's dimensional longitudinal velocity readout in a frame
\end{definition}
\begin{proof}
  This is the declared model definition of particle Longitudinal Velocity In Meters Per Second.
\end{proof}
\begin{definition}[Matches Supplied Velocity Graph]
  \label{def:physics:phyx-mini-0597:phyxminiproblems-problemphyxmini0597-matchessuppliedvelocitygraph}
  \lean{PhyXMiniProblems.ProblemPhyXMini0597.MatchesSuppliedVelocityGraph}
  \uses{def:physics:phyx-mini-0597:phyxminiproblems-problemphyxmini0597-relativisticvelocitygraphsetup, def:physics:phyx-mini-0597:phyxminiproblems-problemphyxmini0597-particlelongitudinalvelocityinmeterspersecond}
  Quantitative and qualitative evidence read from the supplied graph. The curve-to-measurement equality identifies what the graph plots; it does not assign any unshown limiting value to the curve
\end{definition}
\begin{proof}
  This is the declared model definition of Matches Supplied Velocity Graph.
\end{proof}
\begin{definition}[Has Physical Speed Of Light]
  \label{def:physics:phyx-mini-0597:phyxminiproblems-problemphyxmini0597-hasphysicalspeedoflight}
  \lean{PhyXMiniProblems.ProblemPhyXMini0597.HasPhysicalSpeedOfLight}
  \uses{def:physics:phyx-mini-0597:phyxminiproblems-problemphyxmini0597-relativisticvelocitygraphsetup}
  Physical positivity of the dimensional vacuum speed of light
\end{definition}
\begin{proof}
  This is the declared model definition of Has Physical Speed Of Light.
\end{proof}
\begin{definition}[Satisfies One Dimensional Lorentz Velocity Law]
  \label{def:physics:phyx-mini-0597:phyxminiproblems-problemphyxmini0597-satisfiesonedimensionallorentzvelocitylaw}
  \lean{PhyXMiniProblems.ProblemPhyXMini0597.SatisfiesOneDimensionalLorentzVelocityLaw}
  \uses{def:physics:phyx-mini-0597:phyxminiproblems-problemphyxmini0597-relativisticvelocitygraphsetup}
  The governing special-relativistic frame-change law. For every subluminal relative speed v, the independently stored four-velocity in S'(v) is the Lorentz boost of the four-velocity in S by v/c. The desired v → c limit is not included in this law
\end{definition}
\begin{proof}
  This is the declared model definition of Satisfies One Dimensional Lorentz Velocity Law.
\end{proof}
\begin{definition}[Answer Choice]
  \label{def:physics:phyx-mini-0597:phyxminiproblems-problemphyxmini0597-answerchoice}
  \lean{PhyXMiniProblems.ProblemPhyXMini0597.AnswerChoice}
  Labels of the four answer choices printed with the problem
\end{definition}
\begin{proof}
  This is the declared model definition of Answer Choice.
\end{proof}
\begin{definition}[displayed Answer Velocity In Meters Per Second]
  \label{def:physics:phyx-mini-0597:phyxminiproblems-problemphyxmini0597-displayedanswervelocityinmeterspersecond}
  \lean{PhyXMiniProblems.ProblemPhyXMini0597.displayedAnswerVelocityInMetersPerSecond}
  \uses{def:physics:phyx-mini-0597:phyxminiproblems-problemphyxmini0597-answerchoice}
  The dimensional velocity represented by each displayed answer choice
\end{definition}
\begin{proof}
  This is the declared model definition of displayed Answer Velocity In Meters Per Second.
\end{proof}
\begin{definition}[recorded Dataset Answer]
  \label{def:physics:phyx-mini-0597:phyxminiproblems-problemphyxmini0597-recordeddatasetanswer}
  \lean{PhyXMiniProblems.ProblemPhyXMini0597.recordedDatasetAnswer}
  \uses{def:physics:phyx-mini-0597:phyxminiproblems-problemphyxmini0597-answerchoice}
  The source dataset records choice C; this metadata is not a premise
\end{definition}
\begin{proof}
  This is the declared model definition of recorded Dataset Answer.
\end{proof}
\begin{lemma}[laboratory velocity eq point eight c]
  \label{lem:physics:phyx-mini-0597:phyxminiproblems-problemphyxmini0597-laboratory-velocity-eq-point-eight-c}
  \lean{PhyXMiniProblems.ProblemPhyXMini0597.laboratory_velocity_eq_point_eight_c}
  \uses{def:physics:phyx-mini-0597:phyxminiproblems-problemphyxmini0597-relativisticvelocitygraphsetup, def:physics:phyx-mini-0597:phyxminiproblems-problemphyxmini0597-particlelongitudinalvelocityinmeterspersecond, def:physics:phyx-mini-0597:phyxminiproblems-problemphyxmini0597-matchessuppliedvelocitygraph, def:physics:phyx-mini-0597:phyxminiproblems-problemphyxmini0597-hasphysicalspeedoflight, def:physics:phyx-mini-0597:phyxminiproblems-problemphyxmini0597-satisfiesonedimensionallorentzvelocitylaw}
  At zero relative frame speed, the Lorentz law identifies S'(0) with S. Together with the graph's starting ordinate, this determines the particle's laboratory-frame longitudinal velocity as 0.800 c
\end{lemma}
\begin{proof}
  The conclusion follows from the governing relations and figure readouts named in the statement of laboratory velocity eq point eight c.
\end{proof}
% --- Archon declaration topology end ---
% --- Archon physics formalization source end ---
